\documentclass[10pt,a4paper]{amsart}
\usepackage[margin=19mm]{geometry}
\usepackage{amsmath,amssymb,mathtools}
\usepackage[scr=rsfs,cal=euler]{mathalfa}
\usepackage{xfrac}
\usepackage[unicode,hidelinks,bookmarks=false]{hyperref}
\newcommand{\ie}{i.e.\ }
\newcommand{\cf}{cf.\ }
\newcommand{\resp}{respectively\ }
\newcommand{\Subsection}{\subsection}
\providecommand{\nobreakdash}{\nobreak\hbox{-}}
\setlength{\emergencystretch}{2em}
\multlinegap0pt
\usepackage{bm}
\usepackage{bbm}
\usepackage{dsfont}
\usepackage{xspace}
\usepackage{cancel}
\usepackage{mathtools}
\usepackage{cleveref}
\usepackage{amsthm}
\makeatletter
\@ifundefined{lemma}{\newtheorem{lemma}{Lemma}}{}
\@ifundefined{proposition}{\newtheorem{proposition}{Proposition}}{}
\makeatother
\newcommand{\N}{\mathbb{N}}
\newcommand{\R}{\mathbb{R}}
\newcommand{\vd}{\mathrm{d}}
\newcommand{\vx}{\mathbf{x}}
\title[Beyond Mixtures and Products for Ensemble Aggregation: A Lik]{Beyond Mixtures and Products for Ensemble Aggregation: A Likelihood Perspective on Generalized Means}
\author{Rapha\"el Razafindralambo, R\'emy Sun, Fr\'ed\'eric Precioso, Damien Garreau, Pierre-Alexandre Mattei}
\date{Source: arXiv:2603.04204v1 (CC BY 4.0)}
\begin{document}
\maketitle

\noindent\emph{Source and licence.} Beyond Mixtures and Products for Ensemble Aggregation: A Likelihood Perspective on Generalized Means. Version arXiv:2603.04204v1 (CC BY 4.0), distributed under the Creative Commons Attribution 4.0 International licence, as recorded in the arXiv licence metadata for this exact version and shown on the abstract page https://arxiv.org/abs/2603.04204v1. Full rights notice: licenses/arxiv-2603.04204-CC-BY-4.0.txt.\par\smallskip

\noindent\textit{Editorial scope.} Counted unit: the Proposition whose source label is \texttt{proposition: z for 1 over n} in the article (source0.tex lines 999--1038, source label \texttt{proposition: z for 1 over n}), delivered together with the article's own complete original proof of it (source0.tex lines 1041--1072, one \texttt{proof} environment of 32 lines). No mathematics is written, repaired or re-derived: statement and proof are the article's own text line for line. Required context retained uncounted: lemma: multinomial (lines 981--997). Every definition, display and label of those lines is retained unchanged. Omitted from this package and not redistributed: 1--980, 998--998, 1039--1040, 1073--1106. Nothing inside a retained excerpt is omitted or altered: every retained body is delivered exactly as the article's lines are written, so no omission receipt of any kind accompanies this package. Citations: the retained excerpts carry no citation command at all, so no bibliography entry of the article is transcribed and no reference is redistributed. Reference adaptation: none was needed and none was made; every reference inside the delivered unit resolves inside it, so no unresolved reference or citation is produced. Printed slips kept literal: no printed word, symbol, hypothesis or number of the counted statement or of its proof was altered. Layout and class: the article's own class is \texttt{uai2026} with packages including babel, natbib, mathtools, amsthm, amssymb, dblfloatfix, booktabs, tikz, cleveref, todonotes, subcaption, comment, bm; the wrapper is the lane's pinned \texttt{amsart} editorial wrapper at 10pt on A4, which restates the theorem-like environments the retained text needs and adds the article's own macro definitions that the retained text needs (\texttt{\textbackslash N}, \texttt{\textbackslash R}, \texttt{\textbackslash vd}, \texttt{\textbackslash vx}), with extra wrapper packages (bm, bbm, dsfont, xspace, cancel, mathtools, cleveref). The delivered numbering is the unit's own. Assets: no retained span contains a figure environment, a graphics-inclusion command, a PDF-page inclusion, a table or a TikZ picture; no image file is copied into this package, the package declares no asset dependency and no image file is redistributed with it. Licence: version arXiv:2603.04204v1 of this article is distributed under the Creative Commons Attribution 4.0 International licence, as recorded in the arXiv licence metadata for this exact version and shown on the abstract page https://arxiv.org/abs/2603.04204v1. A full rights notice is delivered at licenses/arxiv-2603.04204-CC-BY-4.0.txt. Characters: every retained excerpt is pure ASCII, so the delivered engine, XeLaTeX with its default Latin Modern text font, reports no missing character.
\begin{lemma}[\textbf{Multinomial theorem}]\label{lemma: multinomial}
Let $k$ be a positive integer and $n$ be a non-negative one. 
Let $a_1,\dots,a_k \in F$ where $F$ is a field (e.g., $F=\R$). Then,
\begin{equation}\label{eq: multinomial}
(a_1 + \cdots + a_k)^n
=
\sum_{\substack{n_1+\cdots+n_k=n \\ n_i \ge 0}}
\binom{n}{n_1,\dots,n_k}
\prod_{i=1}^{k} a_i^{\,n_i},
\end{equation}
where the multinomial coefficient is defined as
\begin{equation}
\binom{n}{n_1,\dots,n_k}
=
\frac{n!}{n_1!\cdots n_k!}.
\end{equation}
\end{lemma}

\begin{proposition}\label{proposition: z for 1 over n} Let $p^{(1)},\dots,p^{(k)}$ be $k \geq 1$ Gaussian densities of parameters $(\mu_i,\Sigma_i)$. If $r = \frac{1}{n}$ with $n \in \N \setminus \{0\}$, we have
\begin{equation}
\begin{aligned}
Z_{k,r}
&= \frac{1}{k^n}\sum_{\substack{n_1+\cdots+n_k=n \\ n_i \ge 0}} 
\frac{n!}{n_1!\cdots n_k!}
(2 \pi)^{\frac{d}{2}}
\Big| \sum_{i=1}^k \frac{n_i}{n} \Sigma_i^{-1} \Big|^{-\frac{1}{2}}
\prod_{i=1}^k (2\pi)^{-\frac{d n_i}{2n}}
| \Sigma_i |^{-\frac{n_i}{2n}} \\
&\quad \cdot
\exp\Bigg(
\frac{1}{2}
\Big(\sum_{i=1}^k \frac{n_i}{n} \mu_i^\top \Sigma_i^{-1}\Big)
\Big(\sum_{i=1}^k \frac{n_i}{n} \Sigma_i^{-1}\Big)
\Big(\sum_{i=1}^k \frac{n_i}{n} \Sigma_i^{-1} \mu_i\Big)
-\frac{1}{2}
\sum_{i=1}^k \frac{n_i}{n} \mu_i^\top \Sigma_i^{-1} \mu_i
\Bigg).
\end{aligned}
\end{equation}
If $r=0$, 
\begin{equation}
\begin{aligned}
Z_{k,r} = & (2 \pi)^{\frac{d}{2}}
\Big| \sum_{i=1}^k \frac{1}{k} \Sigma_i^{-1} \Big|^{-\frac{1}{2}}
\prod_{i=1}^k (2\pi)^{-\frac{d}{2k}}
| \Sigma_i |^{-\frac{1}{2k}} \\
&\quad \cdot \exp\Bigg(
\frac{1}{2}
\Big(\sum_{i=1}^k \frac{1}{k} \mu_i^\top \Sigma_i^{-1}\Big)
\Big(\sum_{i=1}^k \frac{1}{k} \Sigma_i^{-1}\Big)
\Big(\sum_{i=1}^k \frac{1}{k} \Sigma_i^{-1} \mu_i\Big)
-\frac{1}{2}
\sum_{i=1}^k \frac{1}{k} \mu_i^\top \Sigma_i^{-1} \mu_i
\Bigg).
\end{aligned}
\end{equation}

\end{proposition}

\begin{proof}
From \Cref{lemma: multinomial}, we have
\begin{equation}\label{eq: from lemma d}
Z_{k,r} = \frac{1}{k^n}\sum_{\substack{n_1+\cdots+n_k=n \\ n_i \ge 0}} 
\binom{n}{n_1,\dots,n_k}
\underbrace{\int \prod_{i=1}^k p_i(\vx)^{\frac{n_i}{n}}}_{I_k}.
\end{equation}
We remark that $I_k$ is the integral of a weighted geometric mean since $\sum_{i=1}^k \frac{n_i}{n} = 1$. Let us denote the weights as $w_i = \frac{n_i}{n}$ for all $i$. For all 
$\vx \in \mathcal{X}$, we have
\begin{align}
\prod_{i=1}^k p_i(\vx)^{w_i} & = \prod_{i=1}^k (2\pi)^{-\frac{dw_i}{2}}| \Sigma_i |^{-\frac{w_i}{2}} \exp \left(-\frac{1}{2} \sum_{i=1}^k w_i (\vx - \mu_i)^\top \Sigma_i^{-1} (\vx - \mu_i)\right) \label{eq: weighted geometric mean first expression} \\
& = C_k\exp \left( -\frac{1}{2} \big( \vx^\top Q_k \vx - 2 b_k^\top \vx + c_k \big)\right)
\end{align}
where
\begin{equation}\label{eq: intermediate quantities}
C_k = \prod_{i=1}^k (2\pi)^{-\frac{dw_i}{2}}| \Sigma_i |^{-\frac{w_i}{2}}, \quad Q_k = \sum_{i=1}^k w_i \Sigma_i^{-1} \succ 0, \quad b_k = \sum_{i=1}^k w_i \Sigma_i^{-1} \mu_i, \quad c_k = \sum_{i=1}^k w_i \mu_i^\top \Sigma_i^{-1} \mu_i.
\end{equation}
Completing the square allows us to factor out a term with an explicit integral. Indeed,
\begin{align}
\prod_{i=1}^k p_i(x)^{w_i} & =  C_k \exp \left( -\frac{1}{2} \big( \vx^\top Q_k \vx - 2 b_k^\top \vx \big) \right) \exp\left( -\frac{1}{2} c_k \right)\\
& =  C_k \exp \left( -\frac{1}{2} \big( \vx^\top Q_k \vx - 2 (Q_k^{-1} b_k )^\top Q_k \vx \big) \right) \exp\left( -\frac{1}{2} c_k \right) \label{eq: Q symmetric} \\
& =  C_k \exp \left( -\frac{1}{2} \big( \vx^\top Q_k \vx - 2 (Q_k^{-1} b_k )^\top Q_k \vx \big) + \underbrace{(Q_k^{-1} b_k)^\top Q_k (Q_k^{-1} b_k)}_{b_k^\top Q_k b_k} \right) \exp\left( \frac{1}{2} b_k^\top Q_k b_k -\frac{1}{2} c_k \right)\\& = C_k \exp \left( -\frac{1}{2} \big( \vx - Q_k^{-1}b_k)^\top Q_k (\vx - Q_k^{-1}b_k)\big)\right) \exp\left( \frac{1}{2} b_k^\top Q_k b_k -\frac{1}{2} c_k \right)\\
& = C_k (2 \pi)^{\frac{d}{2}} | Q_k |^{-\frac{1}{2}} \,\, \mathcal{N}(\vx; \, Q_k^{-1}b_k, Q_k^{-1}) \, \exp\left( \frac{1}{2} b_k^\top Q_k b_k -\frac{1}{2} c_k \right)
\end{align}
where \Cref{eq: Q symmetric} follows from the fact that $Q_k$ is invertible and symmetric. Using $\int\mathcal{N}(\vx; \, Q_k^{-1}b_k, Q_k^{-1}) \vd \vx =1 $, we have
\begin{equation}\label{eq: integral of product}
\int \prod_{i=1}^k p_i(\vx)^{w_i} \vd \vx = (2 \pi)^{\frac{d}{2}} | Q_k |^{-\frac{1}{2}}\prod_{i=1}^k (2\pi)^{-\frac{dw_i}{2}}| \Sigma_i |^{-\frac{w_i}{2}} \exp\left( \frac{1}{2} b_k^\top Q_k b_k -\frac{1}{2} c_k \right)
\end{equation}
where $Q_k,b_k,c_k$ are defined in \Cref{eq: intermediate quantities}. We finally use \Cref{eq: from lemma d} to prove the first result of \Cref{proposition: z for 1 over n}.

For the case $r=0$, we just apply the same reasoning from \Cref{eq: weighted geometric mean first expression} leading to \Cref{eq: integral of product} with $w_i = \frac{1}{k}$.
\end{proof}

\end{document}
